% ===== Section 2: phobias, panic and social anxiety (~27 min, break in the middle) =====
\partaccent{mint}
\topicopener{microphone}{Ֆոբիաներ}{Ոմանք կնախընտրեին վայրի կատվի հետ կռվել, քան կարաոկե երգել}{Ֆոբիաներ, սոցիալական տագնապ ու դրանց բազմաթիվ դրսևորումները}{scared}

\begin{frame}{Կարաոկե. ոմանց համար՝ հաճույք}
  \begin{columns}[c]
    \begin{column}{0.46\textwidth}
      \begin{poll}
        \emoji{microphone}\ Բարձրացրու ձեռքդ, եթե հենց հիմա պատրաստ ես երգել կարաոկե բոլորիս առաջ։
      \end{poll}
      \onslide<2->{%
      \begin{itemize}
        \item Ոմանք սիրում են երգել (հաճախ ոչ սթափ) անծանոթների առաջ
        \item Տաղանդների շոուներում՝ միլիոնների առաջ
      \end{itemize}}
    \end{column}
    \begin{column}{0.48\textwidth}
      \centering
      \onslide<3->{%
      \begin{tikzpicture}
        \draw[-{Stealth[length=2.6mm]}, plum, line width=1pt] (0,0) -- (5.2,0) node[below left, tag, font=\small] {ոգևորություն};
        \draw[-{Stealth[length=2.6mm]}, plum, line width=1pt] (0,0) -- (0,3.2) node[above right, tag, font=\small, anchor=west] {երգելու հմտություն};
        \draw[amRed, line width=2pt] (0.3,2.8) .. controls (1.6,1.4) and (3.0,0.6) .. (4.9,0.35);
        \node at (0.75,2.65) {\bigemoji[0.6cm]{nerd}};
        \node at (4.7,0.85) {\bigemoji[0.6cm]{zany}};
        \node[tag, text width=5cm, font=\small] at (2.6,-1.25) {Բըրնեթի վարկածը. փորձ չի արվել, բայց ալկոհոլը, հավանաբար, դեր ունի};
      \end{tikzpicture}}
    \end{column}
  \end{columns}
\end{frame}

% For others: a nightmare. Meme: two buttons
\begin{frame}{Իսկ ոմանց համար՝ մղձավանջ}
  \begin{columns}[c]
    \begin{column}{0.56\textwidth}
      Նրանց համար կարաոկեի առաջարկը հնչում է այսպես. «Մերկ ձեռնածություն արա կենդանի նռնակներով՝ բոլոր նախկին զույգերիդ աչքի առաջ»։\par\medskip
      \begin{itemize}
        \item<2-> Դեմքը գունատվում է, մկանները լարվում են, շնչառությունն արագանում է. դասական «կռվիր կամ փախիր»
        \item<3-> Երգի ու մահացու կռվի միջև կընտրեին կռիվը։ Եթե, իհարկե, այնտեղ էլ հանդիսատես չկա
      \end{itemize}
    \end{column}
    \begin{column}{0.38\textwidth}
      \centering
      \onslide<4->{%
      \begin{meme}[4.0cm]{two_buttons.jpg}
        \memecap[font=\blackfont\scriptsize]{0.33}{0.69}{երգել\\ կարաոկե}
        \memecap[font=\blackfont\scriptsize]{0.68}{0.9}{կռվել վայրի\\ կատվի հետ}
        \memecap[font=\blackfont\scriptsize]{0.5}{0.1}{սոցիալական\\ տագնապով ուղեղ}
      \end{meme}}
    \end{column}
  \end{columns}
\end{frame}

% Where is the harm?
\begin{frame}{Իսկ ո՞րն է վնասը}
  \begin{columns}[c]
    \begin{column}{0.6\textwidth}
      \begin{itemize}
        \item Կարաոկեն ռիսկ չունի։ Եթե, իհարկե, հանդիսատեսը ստերոիդներ ընդունող երաժշտասերներ չեն
        \item<2-> Ամենավատ դեպքում մի քանի անծանոթ կմտածեն, որ վատ ես երգում։ Հետո՞ ինչ
        \item<3-> Բայց ուղեղի համար դա \textbf{իրական վնաս է}. ամոթ, շփոթմունք, հրապարակային նվաստացում
        \item<4-> Դրանց \textbf{հավանականությունն} իսկ բավական է, որ հրաժարվենք շատ բաներից
      \end{itemize}
    \end{column}
    \begin{column}{0.34\textwidth}
      \centering
      \bigemoji[1.2cm]{flushed}\par\medskip
      \onslide<3->{\mascot[2.4cm]{sad}}
    \end{column}
  \end{columns}
\end{frame}

% Mundane fears; chat
\begin{frame}{Ավելի առօրեական վախեր}
  \begin{columns}[c]
    \begin{column}{0.58\textwidth}
      Միլիոնավոր մարդիկ ամեն օր անում են սրանք առանց խնդրի։ Ոմանց համար դրանք սարսափ են.\par\smallskip
      \begin{enumerate}
        \item \emoji{telephone}\ հեռախոսով խոսել (Բըրնեթն էլ է խուսափում)
        \item \emoji{receipt}\ վճարել, երբ հետևումդ հերթ է
        \item \emoji{beer}\ ընկերների համար խմիչք պատվիրել ու հիշել՝ ով ինչ էր ուզում
        \item \emoji{megaphone}\ ելույթ ունենալ
        \item \emoji{mirror}\ վարսավիրանոցում սանրվել
      \end{enumerate}
    \end{column}
    \begin{column}{0.36\textwidth}
      \centering
      \mascot[2.2cm]{panic}\par\medskip
      \onslide<2->{%
      \begin{chatbox}
        Գրիր այն համարները, որոնք քո մասին են։
      \end{chatbox}}
    \end{column}
  \end{columns}
\end{frame}

% From social anxiety to social phobia
\begin{frame}{Տագնապից մինչև ֆոբիա}
  \centering
  \begin{tikzpicture}
    \shade[left color=mint, right color=blush, rounded corners=2mm] (0,0) rectangle (12,0.8);
    \draw[plum, line width=0.9pt, rounded corners=2mm] (0,0) rectangle (12,0.8);
    \node[anchor=north west, text width=5cm, align=left, font=\small, text=plum] at (0,-0.15)
      {\textbf{Սոցիալական տագնապ}\\ մի քիչ ունի գրեթե յուրաքանչյուրը};
    \node[anchor=north east, text width=5.2cm, align=right, font=\small, text=plum, visible on=<2->] at (12,-0.15)
      {\textbf{Սոցիալական ֆոբիա}\\ երբ այն իրականում խանգարում է ապրել};
    \node at (1.0,1.45) {\bigemoji[0.8cm]{sweatsmile}};
    \node[visible on=<2->] at (11.0,1.45) {\bigemoji[0.8cm]{fearful}};
  \end{tikzpicture}\par\vspace{6mm}
  \raggedright
  \onslide<3->{%
  \begin{card}[colback=mint!35!cardwhite]
    Սոցիալական ֆոբիաները ֆոբիաների \textbf{ամենատարածված} տեսակն են։ Որ հասկանանք դրանք, նախ նայենք ֆոբիաներին ընդհանրապես։
  \end{card}}
\end{frame}

% What a phobia is: the spider
\begin{frame}{Ի՞նչ է ֆոբիան}
  \begin{card}[colback=mint!35!cardwhite]
    \centering \textbf{Ֆոբիա}՝ \textbf{անտրամաբանական} վախ ինչ-որ բանի հանդեպ
  \end{card}
  \vspace{1mm}
  \begin{columns}[T]
    \begin{column}{0.47\textwidth}
      \onslide<2->{%
      \begin{card}
        \emoji[1.4em]{spider}\ {\blackfont Սովորական ռեակցիա}\par\smallskip
        Սարդը նստեց ձեռքիդ, ճչացիր ու թափ տվեցիր ձեռքդ։ Հասկանալի է։
      \end{card}}
    \end{column}
    \begin{column}{0.47\textwidth}
      \onslide<3->{%
      \begin{card}[colback=brainpink!30!cardwhite]
        \emoji[1.4em]{spider}\emoji[1.4em]{scream}\ {\blackfont Ֆոբիա}\par\smallskip
        Անվերջ ճչում ես, շուռ ես տալիս սեղանները, ձեռքդ քլորով լվանում, հագուստդ այրում ու մեկ ամիս տնից դուրս չես գալիս։
      \end{card}}
    \end{column}
  \end{columns}
  \vspace{1mm}
  \onslide<4->{\centering Ի վերջո, դա ընդամենը սարդ է։\par}
\end{frame}

% They know it is irrational
\begin{frame}{Գիտեն, որ անտրամաբանական է}
  \begin{columns}[c]
    \begin{column}{0.58\textwidth}
      \begin{itemize}
        \item Ֆոբիա ունեցողները սովորաբար \textbf{լավ գիտեն}, որ վախն անտրամաբանական է
        \item<2-> Սարդերից վախեցողը (arachnophobia) գիտի, որ մետաղադրամի չափ սարդը վտանգավոր չէ
        \item<3-> «Այն քեզ ոչինչ չի անի» ասելը բարի է, բայց \textbf{լրիվ անօգուտ}
        \item<4-> Վախը գիտակցականից ավելի խորն է
      \end{itemize}
    \end{column}
    \begin{column}{0.36\textwidth}
      \centering
      \onslide<4->{%
      \begin{card}[colback=lavender!35!cardwhite]
        \emoji{bell}\ Հիշիր վեբինար 1-ի «արագ ճանապարհը». նշաձև մարմինը կեղևին չի սպասում։
      \end{card}}
      \mascot[2.0cm]{scared}
    \end{column}
  \end{columns}
  \source{Davidson, Soc. Sci. Med., 2005}
\end{frame}

% Simple and complex phobias
\begin{frame}{Պարզ ու բարդ ֆոբիաներ}
  \begin{columns}[T]
    \begin{column}{0.52\textwidth}
      \begin{card}[colback=mint!35!cardwhite]
        {\blackfont Պարզ (սպեցիֆիկ) ֆոբիաներ}՝ կոնկրետ բանի վախ\par\smallskip
        \emoji{knife}\ առարկա՝ դանակներ\par
        \emoji{spider}\emoji{rat}\ կենդանի՝ սարդեր, առնետներ\par
        \emoji{door}\ իրավիճակ՝ վերելակ\par
        \emoji{syringe}\ այլ՝ արյուն, փսխում\par\smallskip
        \onslide<2->{Վերելակով բարձրանալը տևում է վայրկյաններ. եթե, իհարկե, Վիլի Վոնկան չես։}
      \end{card}
    \end{column}
    \begin{column}{0.42\textwidth}
      \onslide<3->{%
      \begin{card}[colback=lavender!35!cardwhite]
        {\blackfont Բարդ ֆոբիաներ}՝ վախ \textbf{համատեքստերից ու իրավիճակներից}\par\smallskip
        \begin{itemize}
          \item ագորաֆոբիա
          \item սոցիալական ֆոբիա
        \end{itemize}
        Դրանց կհասնենք ընդմիջումից հետո։
      \end{card}}
    \end{column}
  \end{columns}
\end{frame}

% Aplysia
\begin{frame}{Aplysia. պարզ կենդանի, հսկա նեյրոններ}
  \begin{columns}[c]
    \begin{column}{0.44\textwidth}
      \centering
      \includegraphics[width=\linewidth]{aplysia_californica.jpg}\par\smallskip
      {\small Ծովային նապաստակ (Aplysia californica)}
    \end{column}
    \begin{column}{0.52\textwidth}
      \begin{itemize}
        \item Ֆոբիայի հիմքում \textbf{ասոցիատիվ ուսուցումն} է (associative learning). ռեակցիան (վախ) կապվում է խթանի (սարդ) հետ
        \item<2-> Դա կարողանում են անգամ ամենապարզ կենդանիները
        \item<3-> 1970-ականներին հենց Aplysia-ի վրա առաջին անգամ տեսան, թե ինչպես են նեյրոնները փոխվում սովորելիս
        \item<4-> Նեյրոններն այնքան մեծ են, որ դրանց մեջ էլեկտրոդ կարելի է մտցնել
      \end{itemize}
    \end{column}
  \end{columns}
  \source{Castellucci \& Kandel, PNAS, 1974}
\end{frame}

% The straw and the Channel Tunnel; correction: cell bodies, not axons
\begin{frame}{Ծղոտն ու Լա Մանշի թունելը}
  \begin{columns}[c]
    \begin{column}{0.48\textwidth}
      \centering
      \begin{tikzpicture}
        \fill[plum!12!cream] (0,0) circle (1.7);
        \draw[plum, line width=1.2pt] (0,0) circle (1.7);
        \fill[peach!60!cardwhite] (-1.5,-0.85) rectangle (1.5,-0.75);
        \node at (0,-0.2) {\bigemoji[1.3cm]{railway}};
        \node[tag, font=\small\bfseries, anchor=south] at (0,1.8) {Լա Մանշի թունել, մոտ 7.6 մ};
        \filldraw[fill=amOrange, draw=plum, line width=0.4pt] (2.4,-1.4) circle (0.06);
        \draw[plum, line width=0.6pt] (2.4,-1.32) -- (2.6,-0.8);
        \node[tag, anchor=south, font=\small] at (2.7,-0.8) {ծղոտ};
        \node[tag, text width=5.4cm] at (0.4,-2.2) {Մասշտաբը պահպանված չէ. իրականում ծղոտն այստեղ անտեսանելի կլիներ};
      \end{tikzpicture}
    \end{column}
    \begin{column}{0.48\textwidth}
      \begin{bookquote}
        Եթե մարդու աքսոնները ծղոտի չափ լինեին, Aplysia-ինը կլինեին Լա Մանշի թունելի չափ։
      \end{bookquote}
      \onslide<2->{%
      \begin{scienceupdate}
        Գրքում՝ մինչև 1 մմ \textbf{աքսոններ}։ Իրականում 1 մմ-ի են հասնում \textbf{բջջամարմինները}. երևում են անզեն աչքով։
      \end{scienceupdate}}
    \end{column}
  \end{columns}
  \only<2->{\source{National Resource for Aplysia, Univ. of Miami}}
\end{frame}

% How a phobia is learned (link to webinar 1's cake)
\begin{frame}{Ինչպես է ծնվում ֆոբիան}
  \begin{columns}[c]
    \begin{column}{0.3\textwidth}
      \begin{card}[colback=lavender!35!cardwhite]
        \centering\bigemoji[0.8cm]{cake}\ \raisebox{0.25cm}{$\rightarrow$}\ \bigemoji[0.8cm]{nauseated}\par\smallskip
        \raggedright Հիշո՞ւմ ես վեբինար 1-ի տորթը. մեկ անգամ վատացար, ու այլևս չես ուզում։ Նույն մեխանիզմն է։
      \end{card}
    \end{column}
    \begin{column}{0.66\textwidth}
      \centering
      \begin{tikzpicture}[st/.style={pbox, text width=7.6cm, align=left}, ar/.style={-{Stealth[length=2.2mm]}, plum, line width=1pt}]
        \node[st] (s1) at (0,0) {\emoji{megaphone}\ Քեզ զգուշացնում են X-ի մասին};
        \node[st, visible on=<2->] (s2) at (0,-1.0) {\emoji{thought}\ Ուղեղը պատկերացնում է վատ սցենարներ};
        \node[st, visible on=<3->] (s3) at (0,-2.0) {\emoji{runner}\ Հանդիպում ես X-ին. «կռվիր կամ փախիր»};
        \node[st, visible on=<4->, fill=brainpink!30!cardwhite] (s4) at (0,-3.0) {\emoji{bell}\ Նշաձև մարմինը հիշողությանը կպցնում է «\textbf{ՎՏԱՆԳ}»};
        \node[st, visible on=<5->] (s5) at (0,-4.0) {\emoji{scream}\ Հաջորդ անգամ՝ նույն վախը};
        \draw[ar, visible on=<2->] (s1) -- (s2); \draw[ar, visible on=<3->] (s2) -- (s3);
        \draw[ar, visible on=<4->] (s3) -- (s4); \draw[ar, visible on=<5->] (s4) -- (s5);
      \end{tikzpicture}
    \end{column}
  \end{columns}
\end{frame}

% Anything can become a phobia
\begin{frame}{Ֆոբիա կարող է դառնալ ամեն ինչ}
  \begin{columns}[c]
    \begin{column}{0.6\textwidth}
      \begin{itemize}
        \item \emoji{cheese}\ \textbf{Տուրոֆոբիա}՝ պանրի վախ
        \item<2-> \emoji{lemon}\ \textbf{Քսանթոֆոբիա}՝ դեղին գույնի վախ (ակնհայտորեն հատվում է նախորդի հետ)
        \item<3-> \emoji{books}\ Երկար բառերի վախ (որովհետև հոգեբանները չար են).\par
          {\blackfont\small hippopotomonstrosesquipedaliophobia}
        \item<4-> \emoji{fearful}\ \textbf{Ֆոբոֆոբիա}՝ ֆոբիա ունենալու վախ
      \end{itemize}
    \end{column}
    \begin{column}{0.34\textwidth}
      \onslide<4->{%
      \begin{card}[colback=lavender!35!cardwhite]
        Ուղեղը տրամաբանությանը. «Լռի՛ր, դու իմ իսկական հայրիկը չես»
      \end{card}}
      \onslide<5->{%
      \begin{card}[colback=mint!35!cardwhite]
        Բայց որոշ ֆոբիաներ շատ ավելի տարածված են։ Ուրեմն այլ գործոններ էլ կան։
      \end{card}}
    \end{column}
  \end{columns}
\end{frame}

% Social learning: monkeys, snakes and flowers
\begin{frame}{Կապիկը տեսնում է, կապիկը վախենում է}
  \begin{columns}[c]
    \begin{column}{0.58\textwidth}
      \begin{itemize}
        \item Մյուսների վախը տեսնելով՝ կապիկները \textbf{արագ սովորեցին} վախենալ օձերից՝ առանց վարժեցման. \textbf{սոցիալական ուսուցում} (social learning)
        \item<2-> Իսկ եթե օձի փոխարեն \textbf{ծաղիկ} է՞. այդ վախը գրեթե չէին սովորում
      \end{itemize}
      \onslide<4->{%
      \begin{scienceupdate}
        Գրքում՝ շիմպանզեներ։ Իրականում ռեզուս կապիկներ էին, իսկ «ծաղկից վախը» ստացվել էր տեսագրությունը մոնտաժելով։
      \end{scienceupdate}}
    \end{column}
    \begin{column}{0.36\textwidth}
      \centering
      \begin{tikzpicture}
        \node at (0,0) {\bigemoji[1.2cm]{snake}};
        \node at (1.5,0) {\bigemoji[0.9cm]{check}};
        \node[visible on=<2->] at (0,-1.6) {\bigemoji[1.2cm]{rose}};
        \node[visible on=<2->] at (1.5,-1.6) {\bigemoji[0.9cm]{cross}};
        \node[tag, font=\small, text width=4cm, visible on=<3->] at (0.75,-2.85) {Օձի վախը հեշտ է փոխանցվում, ծաղկինը՝ ոչ};
      \end{tikzpicture}
    \end{column}
  \end{columns}
  \only<1-3>{\source{Mineka \& Cook, 1988}}%
  \only<4->{\source{Cook \& Mineka, J. Abnorm. Psychol., 1989}}
\end{frame}

% Evolved fears
\begin{frame}{Ինչից ենք «ծրագրավորված» վախենալ}
  \begin{columns}[T]
    \begin{column}{0.48\textwidth}
      \begin{card}
        \emoji[1.3em]{snake}\emoji[1.3em]{spider}\ {\blackfont Օձեր ու սարդեր}\par
        Հնարավոր մահացու վտանգ. էվոլյուցիան մեզ կասկածամիտ է դարձրել
      \end{card}
      \onslide<2->{%
      \begin{card}
        \emoji[1.3em]{door}\ {\blackfont Վերելակներ}\par
        «Թակարդ» են, ու դա ահազանգ է միացնում
      \end{card}}
    \end{column}
    \begin{column}{0.48\textwidth}
      \onslide<3->{%
      \begin{card}
        \emoji[1.3em]{syringe}\ {\blackfont Ներարկում ու ատամնաբույժ}\par
        Ցավ ու ներխուժում մարմնի մեջ
      \end{card}}
      \onslide<4->{%
      \begin{card}[colback=brainpink!25!cardwhite]
        \emoji[1.3em]{rose}\ {\blackfont Իսկ ծաղիկնե՞րը}\par
        Ծաղիկներից (անթոֆոբիա) ոչ ոք չի վախենում. եթե միայն շատ չար ծաղկափոշու ալերգիա չունի
      \end{card}}
    \end{column}
  \end{columns}
  \source{Mallan, Lipp \& Cochrane, Cogn. Emot., 2013}
\end{frame}

% Uncanny valley (schematic, after Mori)
\begin{frame}{Չարագուշակ հովիտ (uncanny valley)}
  \begin{columns}[c]
    \begin{column}{0.6\textwidth}
      \centering
      \begin{tikzpicture}
        \begin{axis}[
            width=8.4cm, height=5.4cm,
            xmin=0, xmax=105, ymin=-60, ymax=110,
            axis x line=bottom, axis y line=left, axis line style={plum, -{Stealth[length=2.4mm]}},
            xtick=\empty, ytick=\empty,
            xlabel={որքան է նման մարդու}, ylabel={համակրանք},
            label style={font=\small, text=plum},
            clip=false,
          ]
          \addplot[draw=none, fill=brainpink!35, visible on=<2->] coordinates {(68,-60) (68,110) (90,110) (90,-60)};
          \addplot[smooth, amBlue, line width=2pt, tension=0.6] coordinates
            {(2,5) (20,18) (40,38) (58,58) (66,52) (74,10) (80,-35) (86,-15) (93,45) (100,95)};
          \node[font=\small, text=plum, anchor=north] at (axis cs:12,2) {ռոբոտ};
          \node[font=\small, text=plum, align=right, anchor=south east, visible on=<3->] at (axis cs:50,64) {գուլպա\\ երկու աչքով};
          \node[font=\small, text=pinkdark, align=right, anchor=east, visible on=<2->] at (axis cs:70,-32) {դիակ, զոմբի,\\ «գրեթե մարդ»};
          \node[font=\small, text=plum, anchor=south east] at (axis cs:99,97) {իրական մարդ};
        \end{axis}
      \end{tikzpicture}
    \end{column}
    \begin{column}{0.36\textwidth}
      \begin{itemize}
        \item Դիակներից զգուշանալը (հիվանդություն, վտանգ) գուցե էվոլյուցիոն է
        \item<2-> «Գրեթե մարդը» թվում է \textbf{անկենդան}. չունի իրական մարդու նուրբ նշանները
        \item<3-> Իսկ գուլպան՝ երկու կպցրած աչքով, պարզապես զվարճալի է
      \end{itemize}
    \end{column}
  \end{columns}
  \source{Mori, 1970 (թարգմ.՝ 2012). Գրաֆիկը սխեմատիկ է}
\end{frame}

% Fears handed down from parents
\begin{frame}{Վախը՝ ժառանգություն}
  \begin{columns}[c]
    \begin{column}{0.58\textwidth}
      \begin{itemize}
        \item «Ավելի լավ է ապահով, քան ափսոս». եթե ուրիշը վախենում է, մենք էլ հավանաբար կվախենանք
        \item<2-> Հատկապես մանկության մեջ, երբ աշխարհը ճանաչում ենք մեծերի միջոցով
        \item<3-> Մայրիկը մուկ տեսնելիս ճչում ու թափահարում է ձեռքերը. երեխայի համար դա վառ ու ցնցող տպավորություն է
      \end{itemize}
    \end{column}
    \begin{column}{0.36\textwidth}
      \centering
      \bigemoji[1.1cm]{mouse}\ \raisebox{0.3cm}{$\rightarrow$}\ \bigemoji[1.1cm]{scream}\par\medskip
      \onslide<4->{%
      \begin{card}[colback=mint!35!cardwhite]
        Ծնողի ֆոբիան հաճախ անցնում է երեխային. ինչպես մի տարօրինակ հին հագուստ, որ փոխանցվում է ավագից կրտսերին։
      \end{card}}
    \end{column}
  \end{columns}
\end{frame}

% Extinction: Pavlov's dogs
\begin{frame}{Պավլովի շներն ու «մարումը»}
  \begin{columns}[c]
    \begin{column}{0.34\textwidth}
      \centering
      \includegraphics[width=\linewidth]{pavlov_dog.jpg}\par\smallskip
      {\small Պավլովի շներից մեկը (թանգարանային նմուշ)}
    \end{column}
    \begin{column}{0.62\textwidth}
      \centering
      \begin{tikzpicture}[row/.style={pbox, text width=7.0cm, align=left, minimum height=0.9cm}]
        \node[row] (r1) at (0,0) {\emoji{bell}\,+\,\emoji{burger}\ $\rightarrow$\ \emoji{drool}\quad զանգ ու սնունդ՝ թուք};
        \node[row, visible on=<2->] (r2) at (0,-1.1) {\emoji{bell}\ $\rightarrow$\ \emoji{drool}\quad զանգն արդեն բավական է};
        \node[row, visible on=<3->, fill=mint!35!cardwhite] (r3) at (0,-2.2) {\emoji{bell}\emoji{bell}\emoji{bell}\ $\rightarrow$\ \emoji{neutral}\quad զանգ առանց սննդի. կապը մարում է};
      \end{tikzpicture}\par\medskip
      \raggedright
      \onslide<4->{Սա \textbf{մարում} է (extinction)՝ ոչ թե դինոզավրերի պես անհետացում։ Ուղեղը սովորում է, որ զանգից ոչինչ չի հետևում։}
    \end{column}
  \end{columns}
  \source{Bouton \& Bolles, Learn. Motiv., 1979}
\end{frame}

% Why phobias do not fade: the circular logic
\begin{frame}{Ֆոբիան ինքն իրեն է «ապացուցում»}
  \begin{columns}[c]
    \begin{column}{0.32\textwidth}
      Ֆոբիայի առարկայի հետ գրեթե ամեն հանդիպում անցնում է \textbf{առանց վնասի}։ Ուրեմն ֆոբիան էլ պիտի մարե՞ր։\par\medskip
      \onslide<6->{%
      \begin{card}[colback=brainpink!30!cardwhite]
        Ոչ. ֆոբիան \textbf{ամրանում է}։
      \end{card}}
    \end{column}
    \begin{column}{0.64\textwidth}
      \centering
      \begin{tikzpicture}[cyc/.style={pbox, text width=3.3cm}, ar/.style={-{Stealth[length=2.6mm]}, plum, line width=1.3pt}]
        \node[cyc, fill=lavender!40!cardwhite, visible on=<2->] (a) at (0,1.95) {«X-ը վտանգավոր է»};
        \node[cyc, fill=peach!45!cardwhite, visible on=<3->] (b) at (2.75,0) {Հանդիպում ես X-ին.\\ «կռվիր կամ փախիր»};
        \node[cyc, fill=brainpink!35!cardwhite, visible on=<4->] (c) at (0,-1.95) {Ադրենալին, դող, խուճապ. տհաճ է};
        \node[cyc, fill=butter!60!cardwhite, visible on=<5->] (d) at (-2.75,0) {«Մարմինս գժվեց. ուրեմն ճիշտ էի»};
        \draw[ar, visible on=<3->] (a.east) to[bend left=25] (b.north);
        \draw[ar, visible on=<4->] (b.south) to[bend left=25] (c.east);
        \draw[ar, visible on=<5->] (c.west) to[bend left=25] (d.south);
        \draw[ar, visible on=<6->] (d.north) to[bend left=25] (a.west);
        \node[visible on=<6->] at (0,0) {\bigemoji[0.9cm]{cyclone}};
      \end{tikzpicture}
    \end{column}
  \end{columns}
\end{frame}

% ---------------------------------------------------------------- break
\begin{frame}[plain]
  \begin{tikzpicture}[remember picture, overlay]
    \fill[butter!60!cream] (current page.south west) rectangle (current page.north east);
    \node at ([yshift=8mm]current page.center) {\mascot[4cm]{sleepy}};
    \node[anchor=north, font=\titlefont\fontsize{30}{32}\selectfont, text=plum] at ([yshift=-17mm]current page.center) {Ընդմիջում};
    \node[anchor=north, font=\mediumfont\large, text=plum] at ([yshift=-31mm]current page.center)
      {5 րոպե \emoji{coffee}\quad Կանգնիր, ձգվիր, ջուր խմիր};
  \end{tikzpicture}
\end{frame}

% Agoraphobia: myth or fact
\begin{frame}{Առասպել, թե՞ փաստ \#1}
  \mythorfact{Ագորաֆոբիան բաց տարածությունների վախ է}{ՄԱՍԱՄԲ}{peachdark}{%
    \begin{card}[colback=peach!35!cardwhite]
      Ագորաֆոբիան վախ է մի տեղից, որտեղից \textbf{հնարավոր չէ փախչել} կամ որտեղ \textbf{օգնություն չկա}։ Ծանր դեպքում մարդը տնից դուրս չի գալիս։
    \end{card}}
  \source{Magee et al., Arch. Gen. Psychiatry, 1996}
\end{frame}

% Panic attacks
\begin{frame}{Խուճապի նոպա}
  \begin{columns}[c]
    \begin{column}{0.5\textwidth}
      \begin{itemize}
        \item Կարող է պատահել \textbf{ցանկացածի} հետ. վախը հաղթահարում է քեզ, ու ոչինչ անել չես կարողանում
        \item<2-> Ախտանիշները մարդուց մարդ տարբեր են
      \end{itemize}
      \vspace{1mm}
      \onslide<2->{%
      \chip[peach]{fearful}{սարսափ} \chip[peach]{lungs}{շնչել չի ստացվում}\par\smallskip
      \chip[peach]{nauseated}{սիրտ է խառնում} \chip[peach]{dizzy}{գլուխը պտտվում է}\par\smallskip
      \chip[peach]{lock}{«թակարդում եմ»}}
    \end{column}
    \begin{column}{0.46\textwidth}
      \onslide<3->{%
      \begin{card}
        {\blackfont\footnotesize \emoji{speech}\ Մի մարդու նկարագրությունը}\par\smallskip
        {\mediumfont «Չեմ կարողանում կանգնել, չեմ կարողանում խոսել։ Ամբողջ մարմնով ուժեղ ցավ եմ զգում, կարծես ինչ-որ բան ինձ սեղմում է փոքրիկ գնդակի մեջ։ Երբ շատ վատ է, չեմ կարողանում շնչել ու փսխում եմ»}
      \end{card}}
    \end{column}
  \end{columns}
  \only<3->{\source{Holmes \& Scheller, Huffington Post, 2014 (գրքից)}}
\end{frame}

% Panic disorder
\begin{frame}{Երբ վախը գալիս է առանց պատճառի}
  \begin{columns}[c]
    \begin{column}{0.58\textwidth}
      \begin{itemize}
        \item Երբեմն ուղեղը շրջանցում է «միջնորդին» ու վախ է առաջացնում \textbf{առանց պատճառի}
        \item<2-> Պատճառ չկա, ուրեմն անելու բան էլ չկա, ու վախը դառնում է ճնշող
        \item<3-> Սա \textbf{խուճապային խանգարում} է (panic disorder)
        \item<4-> Մարդը կապում է անվնաս վայրերը խուճապի հետ ու սկսում վախենալ դրանցից. ահա և ֆոբիա
      \end{itemize}
    \end{column}
    \begin{column}{0.36\textwidth}
      \centering
      \begin{tikzpicture}
        \node[pbox=cardwhite, text width=3.0cm] (n) at (0,1.4) {\emoji{question}\ պատճառ};
        \draw[pinkdark, line width=2pt] (-1.3,1.1) -- (1.3,1.7);
        \node at (0,-0.4) {\mascot[2.3cm]{panic}};
      \end{tikzpicture}
    \end{column}
  \end{columns}
\end{frame}

% Why panic disorder happens: unknown, several theories
\begin{frame}{Ինչու՞. դեռ հայտնի չէ}
  \begin{columns}[T]
    \begin{column}{0.47\textwidth}
      \begin{card}[colback=mint!35!cardwhite]
        \emoji[1.3em]{bandage}\ Նախկին \textbf{տրավմա}, որ ուղեղը դեռ չի «մարսել»
      \end{card}
      \onslide<3->{%
      \begin{card}[colback=peach!40!cardwhite]
        \emoji[1.3em]{dna}\ \textbf{Գենետիկա}. հիվանդի հարազատներն ավելի հաճախ են ունենում
      \end{card}}
    \end{column}
    \begin{column}{0.47\textwidth}
      \onslide<2->{%
      \begin{card}[colback=lavender!35!cardwhite]
        \emoji[1.3em]{testtube}\ Որոշ նեյրոմեդիատորների \textbf{ավելցուկ կամ պակաս}
      \end{card}}
      \onslide<4->{%
      \begin{card}[colback=butter!50!cardwhite]
        \emoji[1.3em]{cyclone}\ \textbf{Աղետալի մտածողություն}. փոքր խնդրից՝ հսկա անհանգստություն
      \end{card}}
    \end{column}
  \end{columns}
  \vspace{2mm}
  \onslide<5->{\centering \emoji{book}\ Բըրնեթ. «Անհիմն վախի հարցում ուղեղը տարբերակների պակաս չունի»\par}
  \source{Knowles et al., 1998; Witvrouw et al., 2009}
\end{frame}

% Social phobia: the approval of others
\begin{frame}{Սոցիալական ֆոբիա}
  \begin{columns}[c]
    \begin{column}{0.56\textwidth}
      \begin{itemize}
        \item Ուրիշների \textbf{բացասական արձագանքի} վախ. օրինակ՝ կարաոկեի հանդիսատեսի
        \item<2-> Թշնամանք պետք չէ. բավական է \textbf{հավանություն չտալը}
        \item<3-> Ուղեղը ուրիշների միջոցով է «չափաբերում» աշխարհն ու իր տեղն այնտեղ
        \item<4-> Իսկ ի՞նչ է փառքը, եթե ոչ անծանոթների հավանությունը
      \end{itemize}
    \end{column}
    \begin{column}{0.38\textwidth}
      \centering
      \bigemoji[0.9cm]{star}\bigemoji[0.9cm]{clapping}\par\smallskip
      \onslide<5->{%
      \begin{bookquote}
        Գուցե բոլոր հայտնի մարդիկ պարզապես զանգվածային հավանության կարո՞տ են։ Մի քիչ տխուր է։ Եթե, իհարկե, այս գիրքը գովածներից չեն։
      \end{bookquote}}
    \end{column}
  \end{columns}
\end{frame}

% The recipe of social anxiety
\begin{frame}{Սոցիալական տագնապի բաղադրատոմսը}
  \centering
  \begin{tikzpicture}[op/.style={font=\blackfont\LARGE, text=plum}]
    \node[pbox=lavender!40!cardwhite, text width=3.6cm, minimum height=1.2cm] at (0,0) {\emoji{crystalball}\ վատ ելքեր կանխատեսելու սովորություն};
    \node[op] at (2.45,0) {+};
    \node[pbox=mint!40!cardwhite, text width=3.6cm, minimum height=1.2cm, visible on=<2->] at (4.9,0) {\emoji{handshake}\ ընդունված ու հավանված լինելու կարիք};
    \node[op, visible on=<3->] at (7.35,0) {=};
    \node[pbox=brainpink!35!cardwhite, text width=2.6cm, minimum height=1.2cm, visible on=<3->] at (9.3,0) {\emoji{fearful}\ սոցիալական տագնապ};
  \end{tikzpicture}\par\vspace{2mm}
  \raggedright
  \begin{columns}[T]
    \begin{column}{0.47\textwidth}
      \onslide<4->{%
      \begin{card}
        \emoji[1.3em]{telephone}\ {\blackfont Հեռախոս}\par
        Դեմքի ու մարմնի ազդանշաններ չկան. վախենում ես վիրավորել կամ ձանձրացնել
      \end{card}}
    \end{column}
    \begin{column}{0.47\textwidth}
      \onslide<5->{%
      \begin{card}
        \emoji[1.3em]{receipt}\ {\blackfont Հերթ դրամարկղի մոտ}\par
        Բոլորին պահում ես, ու բոլորը նայում են, մինչ դու մանր դրամ ես հաշվում
      \end{card}}
    \end{column}
  \end{columns}
  \vspace{1mm}
  \onslide<6->{\centering\small Ըստ էության՝ \textbf{կատարողական տագնապ} (performance anxiety). վախ՝ հանդիսատեսի առաջ սխալվելուց\par}
\end{frame}

% Parenting styles
\begin{frame}{Ծնողների դերը}
  \begin{columns}[T]
    \begin{column}{0.47\textwidth}
      \begin{card}[colback=brainpink!25!cardwhite]
        \emoji[1.3em]{angersymbol}\ {\blackfont Չափից շատ քննադատող ծնողներ}\par\smallskip
        Երեխան մշտապես վախենում է նեղացնել նրանց, անգամ չնչին բանի համար
      \end{card}
    \end{column}
    \begin{column}{0.47\textwidth}
      \onslide<2->{%
      \begin{card}[colback=lavender!35!cardwhite]
        \emoji[1.3em]{hug}\ {\blackfont Չափից շատ պաշտպանող ծնողներ}\par\smallskip
        Երեխան չի տեսնում անգամ փոքր բացասական հետևանքներ։ Մեծանալով՝ դրանք նրան անհամաչափ ուժեղ են ազդում
      \end{card}}
    \end{column}
  \end{columns}
  \vspace{2mm}
  \onslide<3->{%
  \begin{card}[colback=mint!35!cardwhite]
    \emoji{warning}\ Անգամ «անծանոթներից զգուշացիր» անվերջ կրկնելը կարող է այդ վախն անհարկի ուժեղացնել։
  \end{card}}
  \source{Lieb et al., Arch. Gen. Psychiatry, 2000}
\end{frame}

% Avoidance: papering over a mouse hole
\begin{frame}{Խուսափելը օգնում է... այսօր}
  \begin{columns}[c]
    \begin{column}{0.5\textwidth}
      \begin{itemize}
        \item Ֆոբիա ունեցողները հաճախ խուսափում են վախեցնող իրավիճակներից (avoidant behaviour)
        \item<2-> Հոգեկան հանգստի համար լավ է
        \item<3-> Բայց որքան խուսափում ես, այնքան երկար է ֆոբիան մնում ուժեղ ու վառ
      \end{itemize}
    \end{column}
    \begin{column}{0.46\textwidth}
      \centering
      \onslide<4->{%
      \begin{tikzpicture}[scale=0.85]
        \fill[peach!55!cardwhite] (-2.6,-1.4) rectangle (2.6,1.5);
        \foreach \y in {-0.9,-0.2,0.5,1.2}{\draw[peachdark!40, line width=0.6pt] (-2.6,\y) -- (2.6,\y);}
        \foreach \x in {-1.9,-0.6,0.7,2.0}{\draw[peachdark!40, line width=0.6pt] (\x,-0.9) -- (\x,-0.2) (\x+0.6,-0.2) -- (\x+0.6,0.5) (\x,0.5) -- (\x,1.2);}
        \draw[plum, line width=1pt] (-2.6,-1.4) -- (2.6,-1.4);
        \filldraw[fill=cardwhite, draw=plum, line width=0.8pt, rotate around={-6:(0,-0.75)}] (-0.75,-1.38) rectangle (0.75,-0.15);
        \fill[butter, rotate around={-30:(-0.75,-0.2)}] (-1.0,-0.28) rectangle (-0.5,-0.12);
        \fill[butter, rotate around={30:(0.75,-0.25)}] (0.5,-0.33) rectangle (1.0,-0.17);
        \node at (1.3,-1.05) {\bigemoji[0.55cm]{mouse}};
      \end{tikzpicture}\par\smallskip
      \begin{bookquote}
        Կարծես մկան բունը թղթով ծածկես. դրսից ամեն ինչ լավ է, բայց կրծողի խնդիրը մնում է։
      \end{bookquote}}
    \end{column}
  \end{columns}
  \source{Richer, Early Child Dev. Care, 1993}
\end{frame}

% The most common phobia (numbers from the book's own source, Magee et al. 1996)
\begin{frame}{Ամենատարածված ֆոբիան}
  \begin{columns}[c]
    \begin{column}{0.52\textwidth}
      \centering
      \begin{tikzpicture}
        \begin{axis}[
            xbar, bar width=0.6cm, bar shift=0pt,
            width=5.6cm, height=4.3cm,
            xmin=0, xmax=17, ymin=-0.6, ymax=2.6,
            ytick={0,1,2}, yticklabels={ագորաֆոբիա, {պարզ\\ ֆոբիա}, {սոցիալական\\ ֆոբիա}},
            yticklabel style={align=right},
            xtick=\empty, axis x line=none, axis y line*=left, axis line style={plum},
            tick label style={font=\small, text=plum},
            nodes near coords={\pgfmathprintnumber[fixed, precision=1, assume math mode=true]{\pgfplotspointmeta}\%},
            nodes near coords style={font=\blackfont\small, text=plum},
            clip=false,
          ]
          \addplot[fill=amBlue, draw=none] coordinates {(6.7,0)};
          \addplot[fill=amOrange, draw=none] coordinates {(11.3,1)};
          \addplot[fill=amRed, draw=none] coordinates {(13.3,2)};
        \end{axis}
      \end{tikzpicture}\par
      {\small Կյանքում գոնե մեկ անգամ ունեցել են, ԱՄՆ}
    \end{column}
    \begin{column}{0.44\textwidth}
      Պարանոյիկ ուղեղ + ուրիշների հավանության կարիք = վախ, որ մեզ \textbf{անկարող} կհամարեն։\par\medskip
      \onslide<2->{%
      \begin{bookquote}
        Ապացույցը. սա այս եզրակացության \strike{իններորդ} \strike{տասներորդ} \strike{տասնմեկերորդ} {\blackfont քսանութերորդ} տարբերակն է։ Ու դեռ համոզված եմ, որ շատերին դուր չի գա։
      \end{bookquote}}
    \end{column}
  \end{columns}
  \source{Magee et al., Arch. Gen. Psychiatry, 1996 (գրքի աղբյուրը)}
\end{frame}
